\section{Présentation FEMTO-ST}
\label{sec:pres_entreprise}

\subsection{FEMTO-ST}
\label{sec:FEMTO_ST}

L'institut FEMTO-ST (Franche-Comté Electronique Mécanique Thermique et Optique - Sciences et Technologies) est une unité mixte de recherche associée au CNRS (UMR 6174) et rattachée simultanément à :

\begin{itemize}[label=$\bullet$]
    \item l'université Marie-et-Louis-Pasteur (UMLP)
    \item l'Ecole nationale supérieure de mécanique et des microtechniques de Besançon (Supmicrotech)
    \item l'université de technologie de Belfort-Montbéliard (UTBM)
\end{itemize}

Cet institut a été créé le $1^{er}$ janvier 2004 par la fusion de cinq laboratoires de Franche-Comté actifs dans les domaines de la mécanique, de l'optique et des télécommunications, de l'électronique, du temps-fréquence, de l'énergétique et de la fluidique. Au $1^{er}$ janvier 2012, il incorpore le laboratoire d'informatique de l'UMLP, qui devient le $7^{e}$ département de l'institut. \\

FEMTO-ST compte aujourd'hui plus de 700 membres et constitue l'un des plus gros instituts de recherche dans les sciences de l'ingénieur en France. Il est présent en Franche-Comté sur trois sites géographiques, à Besançon, Belfort et Montbéliard.

\subsection{DISC}
\label{sec:DISC}

Le DISC (Département d'Informatique et Systèmes Complexes) est le département de FEMTO-ST qui traite des sciences informatiques. L'objectif principal des recherches menées étant de modéliser, simuler, développer et valider les systèmes complexes. Le département, qui compte environ 100 personnes, est structuré en 4 équipes de recherche :

\begin{itemize}[label=$\bullet$]
    \item Algorithmique Numérique Distribuéee (AND)
    \item Design, Evaluation and Optimisation of Distributed and Intelligent Systems (DEODIS)
    \item Optimisation, Mobility, NetworkIng (OMNI)
    \item Vérification et validation de logiciels et de systèmes embarqués (VESONTIO)
\end{itemize}

Ce stage se déroule au sein de l'équipe VESONTIO dont l'objectif est porté sur la qualité logicielle et le développement de travaux pour améliorer ou analyser la fiabilité, la sécurité ou la performance des systèmes.

\subsection{Projet ADAPT}

Le projet ANR-23-CE25-0004, intitulé ADAPT (Adaptation dynamique de systèmes à composants hiérarchiques), a pour objectif principal de développer un cadre formel et des outils permettant de concevoir des systèmes cyber-physiques adaptatifs de manière correcte par construction. Il s'intéresse en particulier aux systèmes composés d'un grand nombre de composants modulaires, comme les robots modulaires Blinky Blocks, capables de se reconfigurer dynamiquement pour accomplir une tâche commune.

Les objectifs scientifiques du projet sont les suivants :
\begin{enumerate}
    \item Étendre les modèles à composants hiérarchiques afin de permettre la description de systèmes complexes 
    composés de nombreux modules. Cette modélisation hiérarchique doit faciliter la gestion de la complexité en 
    représentant les composants à différents niveaux d'abstraction tout en intégrant les contraintes 
    fonctionnelles et de ressources.
    \item Intégrer des mécanismes de contrôle adaptatif directement dans ces modèles. L'objectif est de permettre 
    aux systèmes de modifier automatiquement leur configuration en fonction des conditions d'exécution, qu'elles 
    soient :
    \begin{itemize}
        \item cyber, comme l'état du système, les interactions entre composants ou la topologie du réseau 
        \item physiques, comme la consommation d'énergie, le bruit des capteurs ou les perturbations de l'environnement.
    \end{itemize}
    \item Implémenter ces mécanismes d'adaptation dans une infrastructure de modélisation fondée sur les systèmes 
    à composants, afin de fournir une chaîne complète allant de la modélisation jusqu'à l'exécution des stratégies 
    d'adaptation.
    \item Valider les mécanismes proposés par simulation avant leur déploiement sur des plateformes robotiques 
    réelles. Les expérimentations sont réalisées à l'aide des outils DR-BIP et VisibleSim, puis appliquées à des 
    robots modulaires tels que les Blinky Blocks.
\end{enumerate}

Au sein de ce projet, il y a le développement d'un langage de programmation dédié à la modélisation formelle de systèmes distribués complexes, appelé \textbf{CHIPS} (Control of Hierarchical Interconnected Programmable Systems). Ce langage est basé sur le langage BIP (Behavior, Interaction, Priority) et permet de décrire des systèmes hiérarchiques composés de composants modulaires. \\

Il veut permettre d'intégrer des mécanismes de contrôle adaptatif directement dans les modèles de systèmes à composants hiérarchiques, afin de permettre aux systèmes de modifier automatiquement leur configuration en fonction des conditions d'exécution. \\

\FloatBarrier

\newpage